% MG05b. Después de la inversión en la fibra de frontera y antes de
% subsec:alpha-completa; requiere tab:alpha-ventana y eq:alpha-telescopia.
\subsubsection{Contrôles d'orientation, de frontière et de position du registre}
\label{mg05:controles}

La clôture entière permet de distinguer trois dépendances opératoires :
le sens de propagation de la retenue, la condition au bord et la position du
registre dodécaphasique. Leur comparaison utilise la même fenêtre de douze
blocs du tableau~\ref{tab:alpha-ventana}. Les mots
régionaux \(P,E,\Phi\) et le registre \(K\) déjà généré sont maintenus, sauf la
modification identifiée dans chaque cas. La référence commune est
\(\mathcal C_{12}\) avec \(c_{13}=0\), dont la sortie est
\eqref{eq:alpha-ventana-cero}.

La production prospective du registre signé et son inversion \(K^{(r)}=H_4U^{(r)}/4\) précèdent ces contrôles. La normalisation du lemme~\ref{lem:k-hadamard} est conservée, en particulier \(H_4^2=4I_4\) et \(\det H_4=-16\). L'inversion détermine le registre une fois \(U\) produit ; sa sélection opératoire demeure dans la construction qui précède la clôture. Les variations de \(K\) qui suivent ont pour fonction d'examiner la réponse d'un lecteur déjà défini.

\begin{proposition}[Séparation des dépendances de la clôture]
\label{mg05:dependencias}
Pour les données du tableau~\ref{tab:alpha-ventana}, les quatre modifications suivantes ont des effets distincts :
\begin{enumerate}
\item La propagation des quotients de gauche à droite, initiée avec un
quotient nul, fait passer le troisième bloc de \(352\) à \(353\).
\item La condition \(c_{13}=1\), avec l'orientation initiale, change le dernier bloc de \(663\) en \(664\) et conserve les onze blocs précédents.
\item La rotation \(K'_j=K_{j+1}\), avec \(K_{13}=K_1\), modifie
le mot de sortie même en maintenant \(c_{13}=0\).
\item La perturbation \(K'_1=K_1+1=235\), les onze autres coordonnées étant fixes, change le premier bloc de \(007\) en \(006\) et conserve les autres.
\end{enumerate}
\end{proposition}
\begin{proof}
Pour le premier cas, nous définissons le quotient initial \(b_0=0\) et
appliquons, dans l'ordre \(j=1,\ldots,12\),
\[
 \widetilde s_j=P_j+E_j-\Phi_j-K_j+b_{j-1},\qquad
 \widetilde A_j=[\widetilde s_j]_{1000},\qquad
 b_j=\frac{\widetilde s_j-\widetilde A_j}{1000}.
\]
Cette orientation produit, en deux moitiés consécutives,
\begin{align*}
 \widetilde A_{1:6}&=(007,297,353,570,284,800),\\
 \widetilde A_{7:12}&=(995,285,106,471,379,663).
\end{align*}
La différence apparaît déjà dans \(j=3\), où le quotient nécessaire à
la clôture originale provient de la quatrième position.

Pour le second cas, la dernière somme passe de \(663\) à \(664\).
Les deux ont un quotient euclidien nul ; par conséquent, \(c_{12}=0\) dans les
deux cas et la récurrence coïncide de \(j=11\) à \(j=1\).
Le mot devient
\begin{align*}
 A^{(1)}_{1:6}&=(007,297,352,569,283,800),\\
 A^{(1)}_{7:12}&=(997,285,105,472,380,664).
\end{align*}
L'identité \eqref{eq:alpha-telescopia} enregistre exactement la variation unitaire de l'extrémité droite.

La rotation du troisième cas, évaluée avec le même bord nul et la
récurrence originale, donne
\begin{align*}
 A^{\mathrm{rot}}_{1:6}&=(698,699,763,639,119,004),\\
 A^{\mathrm{rot}}_{7:12}&=(559,429,226,119,926,030).
\end{align*}
Dans ce cas, la retenue gauche est \(c_1=-1\) ; l'extrémité droite reste nulle. La conservation des deux extrémités dans \eqref{eq:alpha-telescopia} maintient l'égalité entière lors du changement de position des douze coefficients.

Enfin, la perturbation de \(K_1\) laisse intacte toute la
récurrence de \(j=12\) à \(j=2\). En \(j=1\), la somme
diminue de sept à six, avec un quotient nul dans les deux cas. On obtient
\begin{align*}
 A^{\mathrm{pert}}_{1:6}&=(006,297,352,569,283,800),\\
 A^{\mathrm{pert}}_{7:12}&=(997,285,105,472,380,663).
\end{align*}
\end{proof}

Les quatre contrôles séparent le sens du transport, la fibre de
frontière, la position des coefficients et la réponse locale. Leur
interprétation conserve le domaine fini \(N=12\). Le transport
de la queue et la normalisation de la section complète sont développés
dans la suite, avec les compatibilités exigées par le passage à une
profondeur arbitraire.
